% Avsnitt 4.11, ur lectures/L04/appendix/b_exercises.md.

\section{Övningar}
\label{c4:sec:exercises}

\begin{aside}
\textbf{Så kontrollerar du ditt arbete.} Varje övning nedan som ber dig skriva en VHDL-modul har
en utdelad självkontrollerande testbänk under \file{lectures/L04/exercises/}. Skriv din modul i
dess katalog \file{exercises/<module>/}, med det entitetsnamn och den \textbf{portordning}
övningen anger, och kör den sedan med GHDL \textemdash{} se \secref{c2:sec:testbenches} för de tre
kommandona.
\end{aside}

% ------------------------------------------------------------------------------------------
\exerciseset{Metastabilitetsbegrepp}

\exercise{Metastabilitet med egna ord}{Resonemang}
\label{c4:ex:meta}
Förklara metastabilitet med egna ord. Beskriv i din förklaring:
\begin{itemize}
  \item Vad som händer inuti en D-vippa när dess ingång ändras för nära den aktiva klockflanken.
  \item Varför vippans utgång tillfälligt kan bli kvar mellan en giltig logisk \code{0} och logisk
    \code{1}.
  \item Varför utgången kan ta en oförutsägbar tid på sig att stabilisera sig.
  \item Varför två grindar nedströms kan tolka den ostabiliserade utgången olika: den ena som
    \code{0}, den andra som \code{1}.
\end{itemize}

Förklara varför den oenigheten kan få olika delar av en digital krets att hamna i inkonsistenta
tillstånd.

\exercise{Setup, hold, och varför en asynkron ingång inte kan undvikas}{Resonemang}
\label{c4:ex:setup}
Setup-tid och hold-tid definierar ett fönster kring den aktiva klockflanken under vilket en vippas
ingång måste förbli stabil.

Förklara:
\begin{itemize}
  \item Varför en synkront genererad signal normalt kan konstrueras så att den uppfyller de här
    tidskraven.
  \item Varför en asynkron ingång, till exempel en mekanisk tryckknapp, inte har någon fast
    relation till systemklockan.
  \item Varför en asynkron ingång därför kan ändras under setup- och hold-fönstret oavsett hur den
    omgivande synkrona logiken är konstruerad.
\end{itemize}

Dra slutsatsen varför metastabilitetsrisken från en asynkron ingång måste hanteras snarare än helt
enkelt undvikas.

% ------------------------------------------------------------------------------------------
\exerciseset{Att konstruera en synkroniserare}

\exercise{Kedjan, ritad och förklarad}{Krets}
\label{c4:ex:chain}
Rita en dubbelvippsynkroniserare för en asynkron ingång. Kretsen ska innehålla en asynkron ingång,
två D-vippor kopplade i serie, en gemensam systemklocka, och en synkroniserad utgång tagen från den
andra vippan.

Förklara med egna ord:
\begin{itemize}
  \item Varför den första vippan är direkt exponerad för den asynkrona ingången.
  \item Varför den första vippan kan bli metastabil.
  \item Hur intervallet mellan den första och den andra klockflanken ger det första steget tid att
    stabilisera sig.
  \item Varför den andra vippans utgång har en mycket lägre sannolikhet att vara metastabil.
  \item Varför det är det andra steget, och inte det första, som resten av designen ska använda.
\end{itemize}

En synkroniserare gör inte sannolikheten för metastabilitet exakt noll. Förklara varför den ändå
sänker risken till en acceptabelt låg nivå i de flesta designer.

\note[Tips]Fundera på hur mycket stabiliseringstid den första vippan får innan dess värde samplas
av den andra vippan.

\exercise{Ett tredje steg ”för säkerhets skull”}{Resonemang}
\label{c4:ex:third}
En kollega föreslår att använda tre synkroniserarvippor på varje asynkron ingång i en långsam
design som klockas i några megahertz, bara för säkerhets skull.

Ett tredje steg köper ytterligare en klockperiod för ett metastabilt värde att stabilisera sig på,
och kostar ytterligare en klockcykels latens på varje synkroniserad ingång. Vid några megahertz är
den extra klockperioden en enorm tid räknat i metastabilitetstermer, vilket är precis därför
förslaget låter förnuftigt, och också därför det redan kan vara onödigt.

\xpart{a} Resonera igenom avvägningen åt båda hållen. Under vilka omständigheter faller den ut till
det tredje stegets fördel, och under vilka omständigheter räcker det andra steget redan? Ditt svar
måste hänga på något annat än klockfrekvensen, eftersom frekvensen ensam inte avgör saken: säg vad
det hänger på i stället.

\xpart{b} \Secref{c4:sec:mtbf} kom fram till MTBF utifrån den stabiliseringstid som finns per steg.
Utan att göra om den räkningen, säg åt vilket håll var och en av de här faktorerna drar beslutet,
och varför:
\begin{itemize}
  \item Hur ofta den asynkrona ingången faktiskt ändras.
  \item Metastabilitetsegenskaperna hos den FPGA-familj du riktar in dig på.
  \item Vad ett synkroniseringsfel skulle kosta om det inträffade.
\end{itemize}

\xpart{c} Ge din rekommendation för en vanlig långsam design av det slag som förekommer i den här
boken. Beskriv sedan en design, verklig eller påhittad, där du i stället skulle argumentera för det
tredje steget, och säg vilken av faktorerna i \textbf{b)} som gör jobbet i det fallet.

\begin{aside}
\textbf{Kontrollera att ditt svar täcker} de två saker ett tredje steg ger, ytterligare en
klockperiods stabiliseringstid och en lägre sannolikhet att ett ostabiliserat värde når den
funktionella logiken, och de två det kostar, ytterligare en cykels latens och ytterligare vippor.
Om ditt svar på \textbf{a)} vilar på klockfrekvensen snarare än på den erforderliga medeltiden
mellan fel, läs \secref{c4:sec:mtbf} en gång till innan du går vidare.
\end{aside}

\exercise{Aktivera asynkront, släpp synkront}{Resonemang}
\label{c4:ex:resetpattern}
Förklara resetmönstret \textbf{aktivera asynkront, släpp synkront} från \secref{c4:sec:resetsync}.

Beskriv de två delarna var för sig:
\begin{itemize}
  \item Asynkron aktivering: resetten får verkan omedelbart, och väntar inte på en klockflank.
  \item Synkront släpp: resetten tas bort bara vid en aktiv klockflank, och all berörd logik lämnar
    reset i en kontrollerad relation till klockan.
\end{itemize}

Förklara:
\begin{itemize}
  \item Varför omedelbar aktivering av reset kan vara viktig när klockan är stoppad, när klockan är
    instabil, och när systemet måste gå in i sitt resettillstånd utan fördröjning.
  \item Varför ett släpp av en asynkron reset nära en klockflank kan bryta mot en vippas tidskrav
    för recovery eller removal.
  \item Varför synkronisering av släppet hindrar olika vippor från att lämna reset under olika
    klockcykler.
\end{itemize}

% ------------------------------------------------------------------------------------------
\exerciseset{Att bygga kretsen}

\exercise{\code{led_toggle_sync} för hand}{Krets}
\label{c4:ex:circuit}
Bygg kretsen \code{led_toggle_sync} från \secref{c4:sec:worked} för hand i CircuitVerse.

\xpart{a} Lägg till de yttre portarna: ingångarna \code{clock}, \code{reset_n} (asynkron, aktiv
låg) och \code{button_n} (aktiv låg), samt utgången \code{led}.

\xpart{b} Bygg tvåvippsresetsynkroniseraren:
\begin{itemize}
  \item Koppla de två vipporna i serie.
  \item Aktivera båda stegen asynkront med den råa \code{reset_n}-ingången.
  \item Låt resetten fortplanta sig ut ur kedjan synkront.
  \item Använd den synkroniserade resetutgången, \code{reset_s2_n}, i hela resten av kretsen, och
    använd inte den råa \code{reset_n}-ingången direkt utanför resetsynkroniseraren.
\end{itemize}

\xpart{c} Bygg trevippsknappsynkroniseraren och flankdetektorn:
\begin{itemize}
  \item Använd de två första vipporna som dubbelvippsynkroniserare.
  \item Använd den tredje vippan för att lagra det föregående synkroniserade knappvärdet.
  \item Jämför det nuvarande och det föregående synkroniserade värdet, och detektera ett
    knapptryck när det nuvarande synkroniserade, aktivt låga knappvärdet är \code{0} och det
    föregående synkroniserade värdet är \code{1}.
  \item Använd en AND-grind med det nuvarande värdet inverterat för att generera
    \code{button_edge_s2}. Sambandet är
    $\mathit{button\_edge\_s2} = \mathit{previous} \cdot \mathit{current}'$.
\end{itemize}

\xpart{d} Bygg växlingslogiken för lysdioden: lägg till en vippa som lagrar \code{led_s}, resetta
den med \code{reset_s2_n}, växla den bara när \code{button_edge_s2 = 1}, och koppla utgången
\code{led} till \code{led_s}.

\xpart{e} Simulera hela designen med klockperioden \code{1000 ms}. Aktivera och släpp
\code{reset_n}, tryck ned och släpp \code{button_n}, och observera de synkroniserade knappstegen,
\code{button_edge_s2} och \code{led_s}. Avgör:
\begin{itemize}
  \item Om lysdioden växlar vid den första klockflanken efter det fysiska knapptrycket.
  \item Hur många klockcykler som krävs för att knappvärdet ska fortplanta sig genom
    synkroniseraren och flankdetektorn.
  \item Om den observerade latensen stämmer med det beteende det här kapitlet förutsäger.
\end{itemize}

% ------------------------------------------------------------------------------------------
\exerciseset{Studsfiltrering och tajming}

\exercise{Vad kedjan gör och inte gör}{Resonemang}
\label{c4:ex:debounce}
Du har nu byggt och betraktat kretsen, så det här är en fråga om vad du såg.

Trevippskedjan i Övning~\ref{c4:ex:circuit} gör tre jobb på en gång: den minskar risken för
metastabilitet, den lagrar det föregående synkroniserade knappvärdet, och den detekterar ett
knapptryck. Förklara varför den \textbf{inte} i sig garanterar att knappen är studsfiltrerad.

\xpart{a} Skilj på de tre saker som är lätta att blanda ihop:
\begin{itemize}
  \item Synkronisering: gör om en asynkron ingång till en signal som går att använda i
    klockdomänen.
  \item Flankdetektering: genererar en puls när den samplade signalen ändras.
  \item Studsfiltrering: hindrar flera fysiska övergångar från att tolkas som flera separata tryck.
\end{itemize}
Vilka två av dessa gör din krets faktiskt?

\xpart{b} I din simulering gav ett tryck en enda ren växling. Det är ett starkare påstående än det
ser ut: det håller bara under ett särskilt villkor om var studsen hamnade i förhållande till
klockflankerna. Formulera det villkoret exakt, i termer av vad din krets samplade och vad den
därmed såg. Förklara sedan varför det villkoret är en tillfällighet som beror på din klockperiod
snarare än en egenskap hos din krets. Var konkret med vilken av de två, studsen eller samplingen,
din design faktiskt styr.

\xpart{c} Din simulering kördes med klockperioden \code{1000 ms}; FPGA:n går i \code{50 MHz}. En
tryckknapps kontakter studsar typiskt i storleksordningen 1~ms. Räkna ut konsekvenserna själv:
\begin{itemize}
  \item Hur många gånger samplar en klocka med perioden \code{1000 ms} knappen under den
    millisekunden av studs?
  \item Hur många gånger samplar en klocka på \code{50 MHz} den?
  \item Varje samplad övergång ser ut som en ny flank för din flankdetektor. Förutsäg utifrån de
    två talen vad lysdioden gör på kortet vid ett enda fysiskt tryck, och säg varför exakt samma
    krets uppförde sig oklanderligt i CircuitVerse.
\end{itemize}

\xpart{d} Beskriv en åtgärd av vardera slaget: en hårdvaruåtgärd (till exempel ett RC-lågpassfilter
följt av en Schmitt-triggeringång) och en rent digital åtgärd (till exempel att kräva att den
synkroniserade signalen håller en ny nivå under ett visst minsta antal klockcykler innan den
accepteras, eller en räknare som ignorerar ytterligare övergångar under ett fast intervall efter
ett tryck \textemdash{} vilket är vad \chapref{c7:ch}:s timer skulle ge dig).

\xpart{e} Förklara varför dubbelvippsynkroniseraren måste behållas även när en
studsfiltreringsmetod lagts till. Vilket problem löser inte studsfiltreringen?

% ------------------------------------------------------------------------------------------
\exerciseset{Synkroniserarna som egna moduler}

Det genomarbetade exemplet lägger varje synkroniserare i en egen modul, så att varje senare design
instansierar samma två filer i stället för att härleda dem på nytt. Det är därför både
\chapref{c7:ch}:s \code{walking_led} och \chapref{c8:ch}:s \code{fsm_led} ber dig kopiera in de här
två, och därför \chapref{c8:ch}:s \code{seq_detect_mealy}, som inte har någon knapp, tar
resetsynkroniseraren ensam och lämnar den andra därhän.

De två följande övningarna är du som skriver de modulerna. \Secref{c4:sec:resetsync} och
\ref{c4:sec:buttonsync} ger dig logiken; vad de inte ger dig är entiteten runt den, och, för den
andra, steget från en enda knapp till en vektor av dem. Skriv varje modul själv innan du öppnar det
genomarbetade exemplets kopia, och jämför sedan. Övning~\ref{c4:ex:toggle2} sätter därefter ihop
båda till en fungerande design.

\exercise{\code{reset_sync}}{VHDL}
\label{c4:ex:resetsync}
Skriv resetsynkroniseraren som en modul med namnet \code{reset_sync}.

Entiteten har:

\begin{booktable}{@{}p{6em}p{3.4em}p{5.2em}L@{}}
  \thead{Port} & \thead{Riktn.} & \thead{Typ} & \thead{Beskrivning} \\
  \midrule
  \code{clock} & in & \code{std_logic} & 50 MHz systemklocka. \\
  \code{reset_n} & in & \code{std_logic} & Aktiv låg asynkron reset, rakt från omvärlden. \\
  \code{reset_s2_n} & out & \code{std_logic} & Aktiv låg synkroniserad reset, säker för resten av
    designen att använda. Går låg i samma ögonblick som \code{reset_n} gör det, utan att någon
    klockflank är inblandad, och återgår till hög två stigande flanker efter att \code{reset_n}
    gjort det. \\
\end{booktable}

\xpart{a} Implementera \textbf{aktivera asynkront, släpp synkront} (\secref{c4:sec:resetsync}):
\begin{itemize}
  \item En intern signal för det första steget, plus utgången för det andra.
  \item En process, känslig för både \code{clock} och \code{reset_n}.
  \item Vid \code{reset_n = '0'}, driv båda stegen låga, utanför den klockade grenen.
  \item Vid en stigande klockflank, skifta in en konstant \code{'1'} genom de två stegen.
\end{itemize}

\xpart{b} Värdet som tvingas fram vid reset och värdet som skiftas in på klockan är varandras
motsatser, \code{'0'} och \code{'1'}. Förklara i en mening varför det inte är någon motsägelse,
utifrån vad respektive gren är till för.

\xpart{c} Verifiera designen med dess testbänk. Den kontrollerar mönstrets båda halvor var för sig:
att en låg \code{reset_n} drar \code{reset_s2_n} låg utan någon klockflank däremellan, och att ett
släpp av \code{reset_n} lämnar \code{reset_s2_n} låg tills två stigande flanker har passerat.

\xpart{d} Svara i löpande text:
\begin{itemize}
  \item Utgången heter \code{reset_s2_n} snarare än \code{reset_n_sync}. Vad säger \code{s2}, och
    vad skulle behöva ändras i den här modulen för att det namnet skulle bli fel?
  \item Varje annan modul i den här boken tar \code{reset_s2_n} som sin reset. Den här tar den råa
    \code{reset_n}. Varför måste exakt en modul i en design vara undantaget?
\end{itemize}

\note[Tips]Släppet är den intressanta halvan. Vid reset bryr sig processen inte om klockan alls, så
båda stegen går låga tillsammans; på klockan bryr den sig inte om \code{reset_n}, så \code{'1'}:an
behöver en flank per steg för att nå utgången. Två flankers latens vid släpp är priset för hela
mönstret.

\modulefig[0.62]{lectures/L04/appendix/images/reset_sync.png}

\begin{selfcheck}
\textbf{Självkontroll.} Döp din entitet till \code{reset_sync}, med portarna \code{clock},
\code{reset_n} (in) och \code{reset_s2_n} (out), deklarerade i den ordningen; dess testbänk finns i
\file{exercises/reset_sync/}.
\end{selfcheck}

\exercise{\code{button_sync}}{VHDL}
\label{c4:ex:buttonsync}
Skriv knappsynkroniseraren och flankdetektorn som en modul med namnet \code{button_sync}.

Övning~\ref{c4:ex:toggle2} hanterar två knappar genom att skriva en tvåbitars synkroniserare. Den
här är samma krets med bredden lämnad öppen: en generic avgör hur många knappar den betjänar, så
att \code{led_toggle_sync} kan använda den för en knapp och \code{fsm_led} för två utan att
någondera filen behöver ändras. Se \secref{c4:sec:generics} för syntaxen och för varför den här
modulen har en generic där \code{reset_sync} inte har någon.

Entiteten har en generic:

\begin{booktable}{@{}p{4.6em}p{9.4em}p{4em}L@{}}
  \thead{Generic} & \thead{Typ} & \thead{Förval} & \thead{Beskrivning} \\
  \midrule
  \code{COUNT} & \code{natural range 1 to 3} & \code{1} & Antalet knappar, och därmed bredden på
    båda vektorportarna. \\
\end{booktable}

och de här portarna:

\begin{booktable}{@{}p{7em}p{3.4em}p{11.4em}L@{}}
  \thead{Port} & \thead{Riktn.} & \thead{Typ} & \thead{Beskrivning} \\
  \midrule
  \code{clock} & in & \code{std_logic} & 50 MHz systemklocka. \\
  \code{reset_s2_n} & in & \code{std_logic} & Aktiv låg, \textbf{redan synkroniserad} reset. Den
    här modulen synkroniserar inte sin egen reset; det gjorde modulen i
    Övning~\ref{c4:ex:resetsync}. \\
  \code{button_n} & in & \code{std_logic_vector(COUNT - 1 downto 0)} & Aktivt låga asynkrona
    tryckknappar. \\
  \code{button_edge_s2} & out & \code{std_logic_vector(COUNT - 1 downto 0)} & En encykelspuls vid
    varje knapps fallande flank, två klockflanker efter trycket. Varje bit är oberoende av alla
    andra. \\
\end{booktable}

\xpart{a} Bygg trestegskedjan från \secref{c4:sec:buttonsync}, nu över vektorer i stället för
enskilda bitar:
\begin{itemize}
  \item Tre interna \code{std_logic_vector(COUNT - 1 downto 0)}-signaler, en per steg.
  \item Steg 1 och 2 är dubbelvippsynkroniseraren; steg 3 lagrar det föregående värdet av steg 2.
  \item Vid reset, sätt alla tre stegen till \code{(others => '1')}, inte \code{(others => '0')}.
  \item Härled utgången med en enda konkurrent tilldelning, och lägg märke till att samma uttryck
    fungerar oförändrat på vektorer som på bitar.
\end{itemize}

\textbf{Håll den här i minnet.} De två första vipporna är en synkroniserare i sin helhet och
ingenting annat, och Övning~\ref{c5:ex:sync} låter dig bryta ut dem igen som en egen modul,
generisk både i bredd och i det värde de resettar till. Att skriva dem inline här först är vad som
gör den övningen till en befästning i stället för en ny idé.

\xpart{b} Förklara varför resetvärdet är enbart ettor. Enbart ettor betyder ”släppt”, eftersom
knapparna är aktivt låga, så en reset till enbart nollor startar kedjan med påståendet att varje
knapp redan hålls nedtryckt. Följ vad det kostar:
\begin{itemize}
  \item Med knapparna släppta och en reset till enbart nollor, ger modulen ifrån sig en falsk puls
    när det släppta tillståndet fortplantar sig genom de tre stegen? Gå igenom utgångsuttrycket
    flank för flank innan du svarar, i stället för att gissa.
  \item Nu fallet som faktiskt går sönder. En användare håller en knapp nedtryckt i det ögonblick
    resetten släpps. Vad rapporterar modulen med en reset till enbart ettor, och vad rapporterar
    den med enbart nollor? Vilket av de två är det beteende du vill ha, och varför?
\end{itemize}

\xpart{c} Verifiera designen med dess testbänk. Den instansierar din modul \textbf{två gånger}, en
gång med \code{COUNT = 1} och en gång med \code{COUNT = 2}, och kontrollerar att ett tryck ger
exakt en encykelspuls, att en nedhållen knapp inte ger några ytterligare pulser, att ett släpp inte
ger någon alls, och att ett tryck på den ena knappen i ett par aldrig pulsar den andra.

\xpart{d} Svara i löpande text:
\begin{itemize}
  \item \code{COUNT} är begränsad till \code{1 to 3} snarare än lämnad som en obegränsad
    \code{natural}. Vad köper den begränsningen, givet att ingenting i arkitekturen beror på den
    övre gränsen?
  \item Den här modulen är medvetet inte hopslagen med Övning~\ref{c4:ex:resetsync}:s, trots att en
    design som har en knapp alltid behöver båda. Ange den design som motiverar uppdelningen, och
    säg vad den hade behövt göra i stället.
\end{itemize}

\note[Tips]Skriv arkitekturen för \code{COUNT = 1} i huvudet först, och kontrollera sedan att varje
rad du skrev redan är korrekt för en vektor. \code{and}, \code{not} och aggregatet
\code{(others => '1')} fungerar alla elementvis, så generaliseringen ska inte kosta dig någon extra
kod alls. Gör den det är det värt en extra titt.

\modulefig[0.86]{lectures/L04/appendix/images/button_sync.png}

\begin{selfcheck}
\textbf{Självkontroll.} Döp din entitet till \code{button_sync}, med genericen \code{COUNT}
(\code{natural range 1 to 3}, standardvärde \code{1}), ingångarna \code{clock},
\code{reset_s2_n}, \code{button_n} (\code{std_logic_vector(COUNT - 1 downto 0)}) och utgången
\code{button_edge_s2} (\code{std_logic_vector(COUNT - 1 downto 0)}), deklarerade i den ordningen;
dess testbänk finns i \file{exercises/button_sync/}.
\end{selfcheck}

% ------------------------------------------------------------------------------------------
\exerciseset{Att sätta ihop dem till en design}

\exercise{\code{led_toggle_sync2}}{VHDL}
\label{c4:ex:toggle2}
Gör nu \chapref{c3:ch}:s \code{led_toggle} metastabilitetssäker, i VHDL, som en ny modul med namnet
\code{led_toggle_sync2}.

\chapref{c3:ch}:s \code{led_toggle} växlade två lysdioder från två tryckknappar, men matade in
deras råa, asynkrona signaler rakt i flankdetektorn \textemdash{} logiskt korrekt, men osäkert på
riktig hårdvara. Här bygger du om den med varje asynkron ingång skyddad av det här kapitlets
dubbelvippsynkroniserare.

Entiteten har:

\begin{booktable}{@{}p{5.4em}p{3.4em}p{11.4em}L@{}}
  \thead{Port} & \thead{Riktn.} & \thead{Typ} & \thead{Beskrivning} \\
  \midrule
  \code{clock} & in & \code{std_logic} & 50 MHz systemklocka. \\
  \code{reset_n} & in & \code{std_logic} & Aktiv låg asynkron reset. \\
  \code{button_n} & in & \code{std_logic_vector(1 downto 0)} & Två aktivt låga asynkrona
    tryckknappar. \\
  \code{led} & out & \code{std_logic_vector(1 downto 0)} & Två lysdioder; \code{led(i)} växlar en
    gång per tryck på \code{button_n(i)}. \\
\end{booktable}

\xpart{a} Synkronisera varje asynkron ingång:
\begin{itemize}
  \item Tillämpa tvåvippsresetsynkroniseraren från \secref{c4:sec:resetsync}.
  \item Ge knapparna en dubbelvippsynkroniserare plus ett tredje steg för flankdetektering, som i
    \secref{c4:sec:buttonsync} \textemdash{} nu över en \code{std_logic_vector(1 downto 0)}, så att
    båda knapparna synkroniseras parallellt.
  \item Låt aldrig den råa \code{reset_n} eller \code{button_n} nå växlingslogiken; använd bara den
    synkroniserade resetten och flankpulserna per knapp.
\end{itemize}

\note[Tips]Du kan lägga synkroniseraren inline, eller bryta ut den till en subkomponent som tar en
tvåbitars \code{button_n}, med instansieringssyntaxen från \secref{c2:sec:submodules}. Att lägga
den inline är den kortare vägen, och då är bygget en enda fil:

\begin{shell}
ghdl -a --std=93 led_toggle_sync2.vhd led_toggle_sync2_tb.vhd
\end{shell}

Om du i stället bryter ut den, kopiera in de \file{reset_sync.vhd} och \file{button_sync.vhd} du
skrev till Övning~\ref{c4:ex:resetsync} och \ref{c4:ex:buttonsync} (och instansiera
\code{button_sync} med \code{COUNT = 2}), och namnge dem först på raden \code{ghdl -a}, före din
egen modul, så att GHDL ser dem före den fil som instansierar dem:

\begin{shell}
ghdl -a --std=93 reset_sync.vhd button_sync.vhd \
                 led_toggle_sync2.vhd led_toggle_sync2_tb.vhd
\end{shell}

Se \secref{c2:sub:multifile}. Ingen av modulerna delas ut någonstans i kursrepot, eftersom att
skriva dem är Övning~\ref{c4:ex:resetsync} och \ref{c4:ex:buttonsync}.

\xpart{b} Bygg växlingslogiken:
\begin{itemize}
  \item En vippa för lysdiodstillstånd per lysdiod.
  \item Vid den synkroniserade resetten, släck båda lysdioderna.
  \item Växla \code{led(i)} bara när knapp \code{i}:s synkroniserade puls för fallande flank
    utlöses.
  \item Bekräfta att \code{led(i)} svarar på enbart \code{button_n(i)}, aldrig på den andra
    knappen.
\end{itemize}

\xpart{c} Verifiera designen med dess testbänk. Testbänken trycker på varje knapp i tur och ordning
och kontrollerar att reset släcker båda lysdioderna, att inget tryck betyder ingen växling, att ett
tryck på knapp 0 växlar \textbf{bara} \code{led(0)} när synkroniserarlatensen har förflutit, att
hålla den knappen nedtryckt inte växlar den igen, och att ett tryck på knapp 1 sedan växlar bara
\code{led(1)} och lämnar \code{led(0)} som den var.

Var uppmärksam på \textbf{latensen}: testbänken kontrollerar den \emph{exakta} flanken, inte bara
att lysdioden till slut lägger sig rätt. Ett tryck måste nå växlingslogiken genom två
synkroniserarvippor, så \code{led} är fortfarande oförändrad efter den första och den andra
stigande flanken och växlar vid den tredje. Räkna vipporna ett tryck färdas genom och övertyga dig
om att den tredje flanken är den rätta, och att ett defensivt extra pipelinesteg nu skulle fångas
snarare än tolereras.

\xpart{d} Svara i löpande text, genom att jämföra den här designen med din \code{led_toggle} från
\chapref{c3:ch}:
\begin{itemize}
  \item Vad ändrar synkroniseraren när det gäller att trycka på en riktig, studsande knapp?
  \item Varför var den råa versionen osäker att koppla till en fysisk knapp vid \code{50 MHz}?
  \item Din testbänk driver rena, perfekt tajmade övergångar. Vilket av det här kapitlets två
    problem, metastabilitet eller kontaktstuds, kan en testbänk därför aldrig demonstrera för dig?
\end{itemize}

\modulefig[0.68]{lectures/L04/appendix/images/led_toggle_sync2.png}

\begin{selfcheck}
\textbf{Självkontroll.} Döp din entitet till \code{led_toggle_sync2}, med ingångarna \code{clock},
\code{reset_n}, \code{button_n} (\code{std_logic_vector(1 downto 0)}) och utgången \code{led}
(\code{std_logic_vector(1 downto 0)}), deklarerade i den ordningen; dess testbänk finns i
\file{exercises/led_toggle_sync2/}. Den trycker på varje knapp i tur och ordning och kontrollerar
att \code{led(i)} växlar på \code{button_n(i)}:s fallande flank, att den gör det vid rätt
klockflank snarare än bara lägger sig på rätt värde till slut, och att den aldrig utlöser fel
lysdiod.
\end{selfcheck}
